% Compile twice: lualatex 20001230.tex
\documentclass[11pt]{article}
\usepackage[margin=25mm]{geometry}
\usepackage{fontspec,unicode-math}
\setmainfont{Latin Modern Roman}
\setmathfont{Latin Modern Math}
\usepackage{amsmath,amssymb,mathtools}
\usepackage{xurl,hyperref}
\hypersetup{colorlinks=true,urlcolor=blue,linkcolor=blue}
\setcounter{secnumdepth}{0}
\setlength{\parindent}{0pt}
\setlength{\parskip}{5pt}
\setlength{\emergencystretch}{3em}
\newcommand{\sref}[2]{\hyperlink{source:#1:#2}{[S#2]}}
\newcommand{\cataloguescope}[1]{\paragraph{Source question.}#1}
\newcommand{\floor}[1]{\left\lfloor#1\right\rfloor}
\begin{document}
% UnsolvedMath ID 20001230; source catalogue code AIM-COMPUTATION-0068
\setcounter{section}{20001229}
\section{Exact output-sensitive scaling for decomposable discrete models}\label{20001230}
{\small\sffamily UnsolvedMath ID 20001230 · Computer Science}
\subsection{Definitions and mathematical statement}

\paragraph{Shared notation.}Standard number systems are denoted by $\mathbb N,\mathbb Z,\mathbb Q,\mathbb R,\mathbb C$; index ranges, computational models and specialized notation follow the cited problem source and the definitions below.


For an integer matrix $A$, the fiber at $b$ is $\mathcal F_b=\{x\in\mathbb Z_{\ge0}^d:Ax=b\}$. A Markov basis is a finite set of integer moves $u$ with $Au=0$ such that steps $x\mapsto x\pm u$ staying nonnegative connect every nonempty fiber. A table with ten variables having five levels each has $5^{10}$ cells. Algorithmic scaling depends on the interaction model, margin representation, output format, and sparsity of observed counts, as well as on this total number of cells.
\paragraph{Statement recorded in the supplied source.}
Question 53 (Fienberg). How can we make the algorithms and existing theoretical results apply to problems of sizes that are of practical interest? Scaling up to 5 10 tables is a hopeful goal. \sref{20001230}{2}
\paragraph{Scope qualification.}
The source's performance goal is practical scale, not a specified asymptotic bound or benchmark threshold.
\paragraph{Source status.}
The supplied catalogue labels this record partially solved. Its source-status metadata was checked against the supplied snapshot on 2026-10-11. This import does not independently establish current open status or mathematical correctness.
\subsection{Short English statement}
Can algebraic table methods be made practical for tables with ten variables and five categories per variable?
\subsection{Sources}
\begin{itemize}
\item[\hypertarget{source:20001230:1}{S1}] Ulam AI and the UnsolvedMath Contributors, UnsolvedMath: A Curated Collection of Open Mathematics Problems, record 20001230 (AIM-COMPUTATION-0068). Catalogue curation and metadata: CC BY 4.0. The supplied source question is retained; definitions and the English summary are editorial additions.
\url{https://huggingface.co/datasets/ulamai/UnsolvedMath}.
\item[\hypertarget{source:20001230:2}{S2}] AIM Workshop Problem Lists, original problem source linked by catalogue record AIM-COMPUTATION-0068
\url{https://aimath.org/WWN/compalgstat/compalgstat.pdf}.
\end{itemize}
\paragraph{Preparation and review.}
Original question attributed to the cited source and UnsolvedMath contributors. Codex (GPT-6; exact serving revision not exposed) assisted with definitions, English summaries, source inspection and formatting. No new human mathematical authorship or independent proof verification is asserted. Canonical ID 20001230 and source code AIM-COMPUTATION-0068 checked against the complete supplied registry and existing repository statements; known flagged repetitions and obvious identical conjectures excluded. Full mathematical-equivalence/literature review has not been performed. Statement preparation and final review followed contribution prompts C and D; automated checks verify structure and identity, not open status.
\label{end:20001230}
\end{document}
